\documentclass[hidelinks]{article}
\usepackage[a4paper, total={7in, 10in}]{geometry}
\usepackage[dvipsnames]{xcolor}
\usepackage{amsmath}
\usepackage{tikz}
\usepackage{tkz-euclide}
\usepackage[unicode]{hyperref}
\usepackage[all]{hypcap}
\usepackage{fancyhdr}
\usepackage{tocloft} % Add this package
\usepackage{xcolor}
\definecolor{darkgreen}{RGB}{0,100,0}
\usepackage{listings}
\lstset{
  language=R,
  basicstyle=\ttfamily\small,
  keywordstyle=\color{red},
  commentstyle=\color{darkgreen},
  stringstyle=\color{darkgreen},
  numbers=left,
  numberstyle=\tiny,
  stepnumber=1,
  numbersep=5pt,
  frame=single,
  breaklines=true,
  showstringspaces=false
}
\usepackage{amsmath}
\usepackage{amssymb}
\usepackage{pgfplots}
\pgfplotsset{compat=1.18}
\usepackage{amsmath}
\usepackage{float}
\usepackage{amsthm}
\usepackage[shortlabels]{enumitem}



% Remove dots from table of contents
\renewcommand{\cftdot}{}
\renewcommand{\cftsecdotsep}{\cftnodots}
\renewcommand{\cftsubsecdotsep}{\cftnodots}
\newcommand{\RNum}[1]{\uppercase\expandafter{\romannumeral #1\relax}}


\renewcommand{\contentsname}{Table of Contents}
\usetikzlibrary{angles,calc, decorations.pathreplacing}
\definecolor{carminered}{rgb}{1.0, 0.0, 0.22}
\definecolor{capri}{rgb}{0.0, 0.75, 1.0}
\definecolor{brightlavender}{rgb}{0.75, 0.58, 0.89}

\title{\textbf{STAT 545 HW 2}}
\author{Andres Efren Rocha Jayasinha}
\setcounter{secnumdepth}{0}
\date{Updated: \today} % suppresses the printed date

\begin{document}
\hypersetup{bookmarksnumbered=true,}
\pagecolor{white}
\color{black}
\maketitle
\begin{Large}
\tableofcontents
\end{Large}%
\pagebreak

\section{Assignment 2: Limiting and Stationary Distributions}

\subsection{Exercise 3.7}
A Markov chain has \textit{n} states. If the chain is at state \textit{k}, a coin is flipped, whose
heads probability is \textit{p}. If the coin lands heads, the chain stays at \textit{k}. If the coin
lands tails, the chain moves to a different state uniformly at random. Exhibit the
transition matrix and find the stationary distribution.

\smallskip

\noindent \textcolor{darkgreen}{
\textbf{Stationary Distribution.}
Let $X_0, X_1, \ldots$ be a Markov chain with transition matrix $P$.
A \emph{stationary distribution} is a probability distribution $\pi$
that satisfies
\[
\pi = \pi P.
\]
That is,
\[
\pi_j = \sum_i \pi_i P_{ij}, \quad \text{for all } j.
\]
}

\noindent \textbf{Solution}

\noindent A stationary distribution $\pi$ satisfies
\[
\pi = \pi P,
\]
where $\pi$ is a row vector of dimension $(1 \times n)$ and $P$ is an $(n \times n)$ transition matrix.

\medskip

\noindent Writing this product explicitly,
\[
(\pi_1,\ldots,\pi_n)
\begin{pmatrix}
p_{11} & p_{12} & \cdots & p_{1n} \\
p_{21} & p_{22} & \cdots & p_{2n} \\
\vdots & \vdots & \ddots & \vdots \\
p_{n1} & p_{n2} & \cdots & p_{nn}
\end{pmatrix}
=
(\pi_1,\ldots,\pi_n).
\]

\noindent In this chain,
\[
p_{ij} =
\begin{cases}
p, & i=j,\\[4pt]
\dfrac{1-p}{n-1}, & i\neq j.
\end{cases}
\]

\smallskip

\noindent For each fixed $j$, the $j$th coordinate of $\pi P$ is
\[
\pi_j
= \sum_{i=1}^n \pi_i p_{ij}
= \pi_j p + \sum_{i\neq j} \pi_i \frac{1-p}{n-1}.
\]

\noindent Rearranging,
\[
\pi_j(1-p)
= \frac{1-p}{n-1} \sum_{i\neq j} \pi_i.
\]

\noindent Since $1-p>0$, dividing both sides by $1-p$ gives
\[
\pi_j
= \frac{1}{n-1} \sum_{i\neq j} \pi_i.
\]

\noindent Using $\sum_{i=1}^n \pi_i = 1$, we have $\sum_{i\neq j} \pi_i = 1-\pi_j$, hence
\[
\pi_j = \frac{1-\pi_j}{n-1}.
\]

\noindent Solving,
\[
(n-1)\pi_j = 1-\pi_j
\quad\Rightarrow\quad
n\pi_j = 1
\quad\Rightarrow\quad
\pi_j = \frac{1}{n}.
\]

\noindent Therefore, the stationary distribution is
\[
\boxed{\pi = \left(\frac1n,\ldots,\frac1n\right)}.
\]

\subsection{Exercise 3.8}

\noindent Let
\[
P_1 =
\begin{pmatrix}
\frac14 & \frac34 \\
\frac12 & \frac12
\end{pmatrix},
\qquad
P_2 =
\begin{pmatrix}
\frac15 & \frac45 \\
\frac45 & \frac15
\end{pmatrix}.
\]

\noindent Consider a Markov chain on four states whose transition matrix is given by the block matrix

\[
P =
\begin{pmatrix}
P_1 & 0 \\
0 & P_2
\end{pmatrix}.
\]

\begin{enumerate}[(a)]
\item Does the Markov chain have a unique stationary distribution? If so, find it.
\item Does $\lim_{n\to\infty} P^n$ exist? If so, find it.
\item Does the Markov chain have a limiting distribution? If so, find it.
\end{enumerate}

{\color{darkgreen}
\noindent \textbf{Limiting Distribution.}

\noindent Let $X_0, X_1, \ldots$ be a Markov chain with transition matrix $P$.
A \emph{limiting distribution} for the Markov chain is a probability distribution
$\lambda$ with the property that, for all $i$ and $j$,
\[
\lim_{n \to \infty} P^n_{ij} = \lambda_j.
\]
}

\medskip

{\color{darkgreen}
\noindent \textbf{Limiting Distributions are Stationary Distributions.}

\noindent Assume that $\pi$ is the limiting distribution of a Markov chain with
transition matrix $P$. Then $\pi$ is a stationary distribution.
}

\medskip

{\color{darkgreen}
\noindent \textbf{Positive Matrices and Vectors.}

\noindent A matrix $M$ is said to be \emph{positive} if all entries of $M$ are positive.
In this case, we write $M>0$. Similarly, we write $x>0$ for a vector $x$ whose entries
are all positive.

\smallskip

\noindent \emph{Note.} The notation $M>0$ (or $x>0$) means \emph{entrywise positivity},
that is, every entry is strictly greater than zero. This is stronger than
nonnegativity and does not refer to positive definiteness.
}

\medskip

{\color{darkgreen}
\noindent \textbf{Regular Transition Matrix.}

\noindent A transition matrix $P$ is said to be \emph{regular} if some power of $P$
is positive. That is,
\[
P^n > 0 \quad \text{for some } n \ge 1.
\]
}

\medskip

{\color{darkgreen}
\noindent \textbf{Limit Theorem for Regular Markov Chains.}

\noindent Let $X_0, X_1, \ldots$ be a Markov chain whose transition matrix $P$ is regular.
Then the chain has a limiting distribution, which is the unique, positive
stationary distribution of the chain. That is, there exists a unique probability
vector $\pi>0$ such that
\[
\lim_{n\to\infty} P^n_{ij} = \pi_j,
\quad \text{for all } i,j,
\]
where $\pi$ satisfies
\[
\sum_i \pi_i P_{ij} = \pi_j.
\]

\smallskip

\noindent Equivalently, there exists a positive stochastic matrix $\Pi$ such that
\[
\lim_{n\to\infty} P^n = \Pi,
\]
where $\Pi$ has equal rows with common row $\pi$, and $\pi$ is the unique
probability vector satisfying
\[
\pi P = \pi.
\]
}

\noindent \textbf{Solution}

\smallskip

\noindent Part A:

\smallskip

\noindent A stationary distribution
\[
\pi = (\pi_1,\pi_2,\pi_3,\pi_4)
= \big((\pi_1,\pi_2),(\pi_3,\pi_4)\big)
\]
must satisfy \(\pi=\pi P\). Since
\[
P=
\begin{pmatrix}
P_1 & 0\\
0 & P_2
\end{pmatrix},
\]
this is equivalent to requiring that:
\begin{itemize}
\item \((\pi_1,\pi_2)\) is a stationary distribution for \(P_1\),
\item \((\pi_3,\pi_4)\) is a stationary distribution for \(P_2\),
\item \(\pi_1+\pi_2+\pi_3+\pi_4=1\).
\end{itemize}

\smallskip

\noindent We first find the stationary distribution of \(P_1\). Let \((a,b)\) satisfy
\((a,b)=(a,b)P_1\) with \(a+b=1\). Then
\[
a=\frac14 a+\frac12 b
\quad\Longrightarrow\quad
\frac34 a=\frac12 b
\quad\Longrightarrow\quad
3a=2b.
\]
Together with \(a+b=1\), this yields \(a=\frac25\) and \(b=\frac35\). Hence,
\[
\pi^{(1)}=\left(\frac25,\frac35\right).
\]

\smallskip

\noindent Next, we find the stationary distribution of \(P_2\). By symmetry of the transition
matrix \(P_2\), the stationary distribution is
\[
\pi^{(2)}=\left(\frac12,\frac12\right).
\]

\smallskip

\noindent Therefore, any stationary distribution of \(P\) is a combination of these two
stationary distributions. Let \(\alpha\in[0,1]\) denote the total stationary mass assigned
to states \(\{1,2\}\). Then all stationary distributions are given by
\[
\pi(\alpha)
=
\big(
\alpha\,\pi^{(1)},\,
(1-\alpha)\,\pi^{(2)}
\big)
=
\left(
\alpha\frac25,\,
\alpha\frac35,\,
(1-\alpha)\frac12,\,
(1-\alpha)\frac12
\right),
\qquad 0\le \alpha\le 1.
\]

\smallskip

\noindent Consequently, the Markov chain does \emph{not} have a unique stationary distribution.

\medskip

\noindent Part B:

\noindent Since
\[
P=\begin{pmatrix}
P_1 & 0\\
0 & P_2
\end{pmatrix},
\]
it follows by block-matrix multiplication that for every integer $n\ge 1$,
\[
P^n=\begin{pmatrix}
P_1^n & 0\\
0 & P_2^n
\end{pmatrix}.
\]
Hence, $\lim_{n\to\infty}P^n$ exists if and only if both limits
$\lim_{n\to\infty}P_1^n$ and $\lim_{n\to\infty}P_2^n$ exist, in which case
\[
\lim_{n\to\infty}P^n
=
\begin{pmatrix}
\lim_{n\to\infty}P_1^n & 0\\
0 & \lim_{n\to\infty}P_2^n
\end{pmatrix}.
\]

\medskip

\noindent All entries of $P_1$ are strictly positive, so $P_1$ is a regular transition
matrix. By the Limit Theorem for Regular Markov Chains, $P_1$ has a unique
positive stationary distribution, specifically $\pi^{(1)}=(\frac{3}{5},\frac{2}{5})$ as we showed in Part A,
\[
\implies
\lim_{n\to\infty} P_1^n
=
\begin{pmatrix}
\frac{3}{5} & \frac{2}{5}\\
\frac{3}{5} & \frac{2}{5}
\end{pmatrix}.
\]

\medskip

\noindent Similarly, all entries of $P_2$ are strictly positive, so $P_2$ is also regular.
By the Limit Theorem for Regular Markov Chains, $P_2$ has a unique positive
stationary distribution, specifically $\pi^{(2)}=(\frac{1}{2},\frac{1}{2})$ as we showed in Part A,
\[
\implies
\lim_{n\to\infty} P_2^n
=
\begin{pmatrix}
\frac{1}{2} & \frac{1}{2}\\
\frac{1}{2} & \frac{1}{2}
\end{pmatrix}.
\]

\noindent Combining the two block limits, we conclude that
\[
\lim_{n\to\infty}P^n
=
\begin{pmatrix}
\frac25 & \frac35 & 0 & 0\\[4pt]
\frac25 & \frac35 & 0 & 0\\[4pt]
0 & 0 & \frac12 & \frac12\\[4pt]
0 & 0 & \frac12 & \frac12
\end{pmatrix}
\]
Thus, $\lim_{n\to\infty}P^n$ exists and is given by the matrix above.

\medskip

\noindent Part C:

\noindent Recall that a limiting distribution $\lambda$ must satisfy
\[
\lim_{n\to\infty} P^n_{ij} = \lambda_j
\quad \text{for all } i,j,
\]
which is equivalent to requiring that every row of $P^n$ converges to the same
probability vector.

\noindent From part (b), we computed
\[
\lim_{n\to\infty} P^n
=
\begin{pmatrix}
\frac25 & \frac35 & 0 & 0\\[4pt]
\frac25 & \frac35 & 0 & 0\\[4pt]
0 & 0 & \frac12 & \frac12\\[4pt]
0 & 0 & \frac12 & \frac12
\end{pmatrix}.
\]
The first two rows converge to $(\frac25,\frac35,0,0)$, while the last two rows
converge to $(0,0,\frac12,\frac12)$. Since these limiting rows are not equal,
there is no single probability vector $\lambda$ such that
\[
\lim_{n\to\infty} P^n_{ij} = \lambda_j
\quad \text{for all } i.
\]

\noindent Therefore, the Markov chain does not have a limiting distribution.

\subsection{Exercise 3.10}

A Markov chain has transition matrix $P$ and limiting distribution $\pi$.
Further assume that $\pi$ is the initial distribution of the chain; that is,
the chain is in stationarity. Find the following:
\begin{enumerate}[(a)]
\item $\displaystyle \lim_{n\to\infty} \mathbb{P}(X_n = j \mid X_{n-1} = i)$
\item $\displaystyle \lim_{n\to\infty} \mathbb{P}(X_n = j \mid X_0 = i)$
\item $\displaystyle \lim_{n\to\infty} \mathbb{P}(X_{n+1} = k,\; X_n = j \mid X_0 = i)$
\item $\displaystyle \lim_{n\to\infty} \mathbb{P}(X_0 = j \mid X_n = i)$
\end{enumerate}

\noindent \textbf{Solution}

\noindent Part A:

\noindent For each $n\ge 1$, by time-homogeneity of the Markov chain we have
\[
\mathbb{P}(X_n=j \mid X_{n-1}=i)
=
\mathbb{P}(X_1=j \mid X_0=i)
= P_{ij}.
\]
Since this quantity does not depend on $n$, it follows that
\[
\lim_{n\to\infty}\mathbb{P}(X_n=j \mid X_{n-1}=i)=P_{ij}.
\]

\noindent Part B:

\noindent For each $n\ge 0$,
\[
\mathbb{P}(X_n=j \mid X_0=i) = (P^n)_{ij}.
\]
Since the chain has limiting distribution $\pi$, by definition
\[
\lim_{n\to\infty} (P^n)_{ij} = \pi_j \qquad \text{for all } i,j.
\]
Therefore,
\[
\lim_{n\to\infty}\mathbb{P}(X_n=j \mid X_0=i)
=
\lim_{n\to\infty}(P^n)_{ij}
=
\pi_j.
\]

\noindent Part C:

\noindent For $n\ge 0$,
\[
\mathbb{P}(X_{n+1}=k,\;X_n=j \mid X_0=i)
=
\mathbb{P}(X_n=j \mid X_0=i)\,
\mathbb{P}(X_{n+1}=k \mid X_n=j,\;X_0=i).
\]
By the Markov property,
\[
\mathbb{P}(X_{n+1}=k \mid X_n=j,\;X_0=i)
=
\mathbb{P}(X_{n+1}=k \mid X_n=j)
= P_{jk},
\]
and also
\[
\mathbb{P}(X_n=j \mid X_0=i) = (P^n)_{ij}.
\]
Hence
\[
\mathbb{P}(X_{n+1}=k,\;X_n=j \mid X_0=i)=(P^n)_{ij}\,P_{jk}.
\]
Since the chain has limiting distribution $\pi$, we have
\[
\lim_{n\to\infty}(P^n)_{ij}=\pi_j,
\]
and therefore
\[
\lim_{n\to\infty}\mathbb{P}(X_{n+1}=k,\;X_n=j \mid X_0=i)
=
\pi_j\,P_{jk}.
\]

\noindent Part D:

\noindent By Bayes' rule,

\[
\mathbb{P}(X_0=j \mid X_n=i)
=
\frac{\mathbb{P}(X_n=i \mid X_0=j)\,\mathbb{P}(X_0=j)}{\mathbb{P}(X_n=i)}.
\]
Since the chain is started in stationarity, $\mathbb{P}(X_0=j)=\pi_j$ and
$\mathbb{P}(X_n=i)=\pi_i$ for all $n$. Also,
\[
\mathbb{P}(X_n=i \mid X_0=j)=(P^n)_{ji}.
\]
Therefore,
\[
\mathbb{P}(X_0=j \mid X_n=i)=\frac{(P^n)_{ji}\,\pi_j}{\pi_i}.
\]
Because the chain has limiting distribution $\pi$, we have
\[
\lim_{n\to\infty}(P^n)_{ji}=\pi_i,
\]
and hence

\[
\lim_{n\to\infty}\mathbb{P}(X_0=j \mid X_n=i)
=
\lim_{n\to\infty}\frac{(P^n)_{ji}\,\pi_j}{\pi_i}
=
\frac{\pi_i\,\pi_j}{\pi_i}
=
\pi_j,
\]







\subsection{Exercise 3.14}

The California Air Resources Board warns the public when smog levels are above certain thresholds. Days on which the board issues warnings are called \emph{episode days}. Lin (1981) models the daily sequence of episode and nonepisode days as a Markov chain with transition matrix
\[
P =
\begin{pmatrix}
0.77 & 0.23 \\
0.24 & 0.76
\end{pmatrix},
\]
where the rows and columns correspond to the states \emph{Nonepisode} and \emph{Episode}, respectively.

\begin{enumerate}[(a)]
\item What is the long-term probability that a given day will be an episode day?
\item Over a year's time, about how many days are expected to be episode days?
\end{enumerate}

\noindent \textbf{Solution}

\smallskip

\noindent Part A:

\noindent Let $\pi = (\pi_N, \pi_E)$ denote the stationary distribution of the Markov chain, where
$\pi_N$ is the long-term probability of a nonepisode day and $\pi_E$ is the long-term
probability of an episode day.

\smallskip

\noindent The stationary distribution satisfies
\[
\pi = \pi P,
\qquad
\pi_N + \pi_E = 1.
\]

\noindent Using the transition matrix
\[
P =
\begin{pmatrix}
0.77 & 0.23 \\
0.24 & 0.76
\end{pmatrix},
\]
\noindent we obtain the equations
\[
\pi_E = 0.23\,\pi_N + 0.76\,\pi_E.
\]

\noindent Rewriting,
\[
\pi_E - 0.76\,\pi_E = 0.23\,\pi_N
\quad \Longrightarrow \quad
0.24\,\pi_E = 0.23\,\pi_N.
\]

\noindent Using $\pi_N = 1 - \pi_E$, we have
\[
0.24\,\pi_E = 0.23(1 - \pi_E),
\]
which implies
\[
0.47\,\pi_E = 0.23.
\]

\noindent Thus,
\[
\pi_E = \frac{0.23}{0.47} \approx 0.489.
\]

\noindent Therefore, the long-term probability that a given day will be an episode day is approximately
\[
\boxed{0.489}
\]

\noindent Part B:

\noindent From part (a), the long-term probability that a day is an episode day is
\[
\pi_E = \frac{0.23}{0.47} \approx 0.489.
\]
Over a period of $365$ days, the expected number of episode days is therefore
\[
\mathbb{E}[\text{\# episode days in a year}]
\approx 365 \,\pi_E
= 365\left(\frac{0.23}{0.47}\right)
\approx 365(0.489)
\approx 178.617
\]
So, over a year's time we expect about
\[
\boxed{179 \text{ episode days}}
\]

\subsection{Exercise 3.63}

Hourly wind speeds in a northwestern region of Turkey are modeled by a Markov chain
in Sahin and Sen (2001). Seven wind speed levels are the states of the chain.
The transition matrix is
\[
P =
\begin{pmatrix}
0.756 & 0.113 & 0.129 & 0.002 & 0     & 0     & 0     \\
0.174 & 0.821 & 0.004 & 0.001 & 0     & 0     & 0     \\
0.141 & 0.001 & 0.776 & 0.082 & 0     & 0     & 0     \\
0.003 & 0     & 0.192 & 0.753 & 0.052 & 0     & 0     \\
0     & 0     & 0.002 & 0.227 & 0.735 & 0.036 & 0     \\
0     & 0     & 0     & 0.007 & 0.367 & 0.604 & 0.022 \\
0     & 0     & 0     & 0     & 0.053 & 0.158 & 0.789
\end{pmatrix}
\]

\begin{enumerate}[(a)]
\item Find the limiting distribution by
\begin{enumerate}[(i)]
\item taking high matrix powers, and
\item using the \texttt{stationary} command in the \texttt{utilities.R} file.
\end{enumerate}
How often does the highest wind speed occur? How often does the lowest speed occur?

\item Simulate the chain for $100{,}000$ steps and estimate the proportion of times
that the chain visits each state.
\end{enumerate}

\noindent \textbf{Solution}

\smallskip

\noindent Part A:

\noindent \textbf{Input:}

\begin{lstlisting}[caption={Hourly Wind Speeds of a Region In Turkey}]
## Part A

# Part i

# Transition matrix
P <- matrix(c(
  0.756, 0.113, 0.129, 0.002, 0.000, 0.000, 0.000,
  0.174, 0.821, 0.004, 0.001, 0.000, 0.000, 0.000,
  0.141, 0.001, 0.776, 0.082, 0.000, 0.000, 0.000,
  0.003, 0.000, 0.192, 0.753, 0.052, 0.000, 0.000,
  0.000, 0.000, 0.002, 0.227, 0.735, 0.036, 0.000,
  0.000, 0.000, 0.000, 0.007, 0.367, 0.604, 0.022,
  0.000, 0.000, 0.000, 0.000, 0.053, 0.158, 0.789
), nrow = 7, byrow = TRUE)

# Compute a high power of P
Pn <- P
for (i in 1:1000) {
  Pn <- Pn %*% P
}

# Display the result
round(Pn, 6)

# Extract the limiting distribution (any row)
pi_hat <- Pn[1, ]
pi_hat

# Frequencies of lowest and highest wind speed states
pi_hat[1]   # lowest wind speed
pi_hat[7]   # highest wind speed

# Part ii

source("~/Downloads/utilities.R")

station_dist <- stationary(P)
\end{lstlisting}

\noindent \textbf{Output:}

\noindent The Limiting Distribution is (three decimal places): $\begin{pmatrix}
0.324 & 0.206 & 0.303 & 0.131 & 0.029 & 0.002 & 0.000
\end{pmatrix}$

\noindent The highest wind speed occurs 0\% of the time.

\noindent The lowest wind speed occurs 32.4\% of the time.

\medskip

\noindent Part B:

\noindent \textbf{Input:}

\begin{lstlisting}[caption={Hourly Wind Speeds of a Region In Turkey}]
## Part B

P <- matrix(c(
  0.756, 0.113, 0.129, 0.002, 0.000, 0.000, 0.000,
  0.174, 0.821, 0.004, 0.001, 0.000, 0.000, 0.000,
  0.141, 0.001, 0.776, 0.082, 0.000, 0.000, 0.000,
  0.003, 0.000, 0.192, 0.753, 0.052, 0.000, 0.000,
  0.000, 0.000, 0.002, 0.227, 0.735, 0.036, 0.000,
  0.000, 0.000, 0.000, 0.007, 0.367, 0.604, 0.022,
  0.000, 0.000, 0.000, 0.000, 0.053, 0.158, 0.789
), nrow = 7, byrow = TRUE)

# Compute a high power of P
Pn <- P
for (i in 1:100000) {
  Pn <- Pn %*% P
}

P[1, ]
\end{lstlisting}

\noindent \textbf{Output}:

\noindent Each row of the matrix $P_{n}$ converges to (three decimal placces): $\begin{pmatrix}
0.324 & 0.206 & 0.303 & 0.131 & 0.029 & 0.002 & 0.000
\end{pmatrix}$


\end{document}